% Generated draft: requires mathematical and source review.
\documentclass[11pt]{article}
\usepackage[margin=25mm]{geometry}
\usepackage{fontspec}
\setmainfont{Latin Modern Roman}
\usepackage{amsmath,amssymb,mathtools,unicode-math}
\setmathfont{Latin Modern Math}
\usepackage{enumitem,needspace,xurl,hyperref}
\setlength{\parindent}{0pt}
\setlength{\parskip}{5pt}
\setlength{\emergencystretch}{3em}
\newcommand{\N}{\mathbb{N}}
\newcommand{\Z}{\mathbb{Z}}
\newcommand{\Q}{\mathbb{Q}}
\newcommand{\R}{\mathbb{R}}
\newcommand{\C}{\mathbb{C}}
\newcommand{\F}{\mathbb{F}}
\newcommand{\A}{\mathbb{A}}
\newcommand{\PP}{\mathbb{P}}
\newcommand{\E}{\mathbb{E}}
\newcommand{\eps}{\varepsilon}
\newcommand{\dd}{\,\mathrm{d}}
\begin{document}
\section*{UnsolvedMath ID: 20002331}
\textbf{Source problem code:} catalogue entry
\subsection{Definitions and mathematical statement}
Problem 16 asks to study ample generics for an $\omega$-stable, separably categorical continuous theory $T$, generalizing the discrete model-theoretic setting described in the cited small-index paper. Separably categorical means that $T$ has a unique separable model up to isomorphism. Here $\omega$-stability is understood in its continuous-logic sense, involving separability of type spaces in their type metrics. The original wording does not specify whether ``ample generics'' means ordinary or metric ample generics.

For a separable metric structure $M$, equip $G=\operatorname{Aut}(M)$ with the Polish topology $\tau$ of pointwise convergence and the uniform metric
\[
\partial(g,h)=\sup_{x\in M}d(gx,hx)
\]
in a bounded single-sort presentation, or the corresponding specified metric for a many-sorted presentation. Ordinary ample generics requires a comeager diagonal conjugacy orbit in $G^n$ for every integer $n\geq1$. Ample metric generics instead requires that, for every $n\geq1$, some diagonal conjugacy orbit has comeager closure in the uniform product metric, where comeagerness is measured in $\tau^n$. A metric-generic element is the case $n=1$. These are distinct requirements.

The supplied aid reports a quotient obstruction. Suppose $q:G\to A$ is a surjective group homomorphism to a non-discrete Hausdorff abelian topological group, continuous and open for $\tau$ and continuous from the uniform-metric topology on $G$ to the topology on $A$. Then
\[
\overline{\{hgh^{-1}:h\in G\}}^{\partial}\subseteq q^{-1}(\{q(g)\})
\]
for every $g\in G$. The fiber is closed and nowhere dense in $\tau$, so there is no metric-generic element and hence no ample metric generics.

The proposed example is an infinite-dimensional separable complex Hilbert space $H$ in its Hilbert-space language expanded by a named rank-one orthogonal projection $P$. The aid asserts that $\operatorname{Th}(H,P)$ is $\omega$-stable and separably categorical, and
\[
\operatorname{Aut}(H,P)\cong U(PH)\times U((I-P)H)\cong U(1)\times U(\ell_2).
\]
Using the strong operator topology and the uniform metric on the unit ball, equivalently operator-norm distance, the phase projection to $U(1)$ is continuous, open, surjective, and $1$-Lipschitz. The obstruction therefore rules out a metric-generic automorphism in this example. This supplies a negative answer to the universal assertion that the two model-theoretic hypotheses alone force metric ample generics. It does not classify all such theories. The summary also reports positive results for pure Hilbert space and other metric structures in the theory of Polish topometric groups, without giving their full hypotheses or references.

The extraction also includes Problem 17: seek a model-theoretic understanding of parametric families of probability measures and ways to compare different statistical experiments. One may denote such a family by $(\mu_\theta)_{\theta\in\Theta}$, but the source specifies neither the measurable spaces, parameter structures, logic, nor comparison relation between experiments. This remains a separate, unspecified research question; the automorphism-group obstruction above does not answer it. Its trailing ``Background:'' and page marker $4$ carry no mathematical assumptions.
\subsection{Short English statement}
The main problem seeks a metric analogue of ample-generics results for stable categorical structures. The supplied aid gives a phase-quotient obstruction and a Hilbert-space expansion illustrating that the stated hypotheses do not force metric ample generics. A neighboring question about probability measures is retained separately, with its missing framework flagged.
\subsection{Sources}
\begin{itemize}
\item Local UnsolvedMath registry, permanent ID $20002331$, \texttt{catalogue entry}; \texttt{aim-logic-notes.json}, \texttt{notes[106]}.
\item American Institute of Mathematics, \emph{Model theory of metric structures}; supplied verification calls the PDF \emph{Questions on Model theory for metric structures}. Problems 16 and 17 occur in this merged record: \url{https://aimath.org/WWN/continuouslogic/continuouslogic.pdf}.
\item Source background names Hodkinson, Hodges, Lascar and Shelah and the paper \emph{The small index property for $\omega$-stable, $\omega$-categorical structures and the random graph}. No publication data or other authoritative citation for metric ample generics, the positive examples, or the projection expansion is supplied.
\item Research attempt 1: partial result, problem partially solved at run, rigor label ``proved,'' novelty confidence low. Registry classification partial progress, status partially solved, and statement status exact are supplied labels.
\end{itemize}
\end{document}
